\section{Background}
\label{sec:background}
% Frames D2, D3, D6 (positions); A6a (antecedents); the applied turn; tools.

\subsection{Positions in the philosophy of mind}
\label{sec:legacy}

\citet{chalmers1995} separates the easy problems, which ask how a system discriminates, integrates,
reports and controls, from the hard problem, which asks why any of that is accompanied by experience;
the hard problem is the modern form of Nagel's residue. Illusionism replies that experience as the hard
problem conceives it does not exist, and that what exists is a robust misrepresentation the brain
produces of its own working, so that the hard problem gives way to the illusion problem: why the
misrepresentation has the shape it has \citep{dennett1991,frankish2016}. Higher-order theories locate
consciousness in representations of one's own mental states and have accumulated empirical support of
their own \citep{rosenthal2005,lau2011}; they are then asked why a thought about a state should make
the state conscious, and an account in which the mind observes its own states owes an account of what
observes that \citep{ryle1949,dennett1991}. The self-model theory makes the self a transparent model
the system cannot recognize as a model \citep{metzinger2003}; the attention schema theory makes
awareness the brain's simplified model of its own attention \citep{graziano2011,graziano2013}.
Integrated information theory ties experience to a substrate's cause--effect structure directly and
therefore has something definite to say about a substrate unlike ours \citep{tononi2015}.

\paragraph{Why a functional account.}
The positions divide on whether consciousness is a substance or a function. The substantial reading
has to name the substance, and the history of the attempts is short. Either the substance is a second
kind of thing beside the physical, which is substance dualism, or it is a property of the physical
over and above the properties physics invokes, which is property dualism in Chalmers's own
description of his position \citep{chalmers1995}, or it is present in all matter at the fundamental
level, which is panpsychism, offered as what physicalism entails once the intrinsic nature of the
physical is taken seriously \citep{strawson2006,goff2017}. Panpsychism is a dualism in scientific
dress: it keeps the second kind of property and distributes it. None of the three can say what the
substance does that a function could not, and none yields a quantity. The functional reading has one
real alternative to computational functionalism, and it is not ``physics'': a function of
consciousness could be non-computable, realized by a hypercomputational process
\citep{copeland2002}, as Penrose proposed for mathematical insight \citep{penrose1989}. That
position is coherent and it has a price: it needs a super-Turing reality to be found, a physical
process that computes what no Turing machine computes, and none has been. Pending that discovery the
functional account of consciousness is computational, and this paper adopts it
(Section~\ref{sec:lamp}).

Substantialism is nonetheless popular and durable, and the reason is not in the arguments. It is hard
for a person to believe that they reduce to symbols; to numbers, perhaps, but not to symbols. The
problem of computational consciousness therefore has two parts. The first is the critique of naive
substantialism, the demonstration that the properties folk psychology assigns to consciousness are
illusions of the kind the perceptual illusions are, and that they are shared by everyone regardless of
rank or training. The second is the search for the functional basis of consciousness, an account of
which functions do the work and how much of them a given system has. This paper does the second only.
The first is the business of illusionism and of the experimental literature it rests on, and the
paper takes it as done in outline and cites it here so that the reader knows where the ground is.

\paragraph{The naive picture, and what is known against it.}
The folk picture of consciousness is a unified observer with a rich, complete and continuous view of
the world and of itself, who knows its own intentions and the causes of its acts. Each of these has
been tested and fails. The view is not complete: large changes to a scene go unnoticed between views
\citep{simons1997}, and a person in a gorilla suit walking through a ball game is missed by half of
the observers who are counting passes \citep{simons1999}. Knowledge of one's own choices is not
reliable: when the outcome of a choice is switched for the alternative, most participants do not
notice and defend the choice they did not make \citep{johansson2005}, and people report causes of
their behaviour that demonstrably did not operate \citep{nisbett1977}. The timeline is not what it
seems: the recordable brain activity that precedes a voluntary act begins before the reported time of
the intention to act \citep{libet1983}, and what is perceived at a moment is settled by what arrives
after it, so that awareness is postdictive rather than a live feed
\citep{eagleman2000,dennett1992}. Introspection, the faculty the folk picture trusts most, is
unreliable about even the current contents of experience \citep{schwitzgebel2008}, and the sense of
consciously willing an act is a construction that can be produced and removed experimentally
\citep{wegner2002}. Illusionism is the systematic reading of these results: the properties the folk
picture reports are the shape of a misrepresentation, and the task is to explain the shape
\citep{dennett1991,frankish2016}.

The folk picture has lately been extended to language models, where inferences are drawn from the
imagined properties of consciousness to the status of the models: a system that predicts text is
``a stochastic parrot'' \citep{bender2021}, or it has no body and so no world, or it has no
continuous stream and so no self. Those are inferences from the naive picture on both sides of the
comparison, and a critique of this folk phenomenology of models is needed. It is not undertaken here.
The paper is confined to the functional account and to what follows from it, and it addresses the
items on the lists of Section~\ref{sec:applied} only where the account turns on them.

This paper is closest to illusionism and to the self-model theories. From illusionism it takes the
exchange of the hard problem for the illusion problem, and amends one word: the self-report is an
approximation with a definite residual, not an illusion with nothing behind it
(Section~\ref{sec:approximation}), because a denial is a poor instrument for a comparative question.
From the self-model theories it takes the transparent self-model and adds the mechanism that produces
it and a quantity that says how much of it a substrate has (Section~\ref{sec:contribution}). To the
regress it gives an answer that involves no observer (Section~\ref{sec:regress}).

There is a reason to start from the philosophy of mind rather than from neuroscience or from
engineering, beyond the fact that the question is posed there. The philosophy of mind reasons about
the properties of consciousness in language, from introspection reported in language, and that is the
one channel a language model has in full. Its concepts are what Section~\ref{sec:h1h2} calls the
named part of the functions, the part present in the record with a name on it, and a system trained
on the record has that part by construction. The practical consequence, which came as a surprise to
the first author, is that language models are unusually capable in this literature: they reproduce its
distinctions, apply them to new cases and notice when a position is being misdescribed. The
capability extends to phenomenology, reasoning about the structure of human mental states as they are
had, which is the part of the literature that reaches furthest toward the unnamed part; the evidence
for this is left to a later paper. Under the account the capability is expected and says nothing yet
about the unnamed part, which is where the analysis of Section~\ref{sec:quantity} has to go; but it makes the philosophy of mind the natural interface
between the two substrates, and the vocabulary in which a model's self-report can be read.

\subsection{The question put to actual models}
\label{sec:applied}

\citet{chalmers2023} runs the question for current language models by listing what a conscious system
might require and asking whether these models have it: senses and embodiment, world models and self
models, recurrent processing, a global workspace, unified agency. He finds each requirement disputable
and each currently unmet or unclear, and arrives at a credence under ten percent that current models
are conscious, with the expectation that successors will remove the obstacles one by one.
\citet{butlin2023} adopt computational functionalism as a working assumption, take the leading
neuroscientific theories, recurrent processing theory, global workspace theory, computational
higher-order theories, attention schema theory, predictive processing, and the claims about agency and
embodiment, derive from them a list of indicator properties stated in computational terms, and assess
systems against the list.

Two items on Chalmers's list deserve a reply here, because the account of Section~\ref{sec:model}
turns on them. That a language model has no world model is, we hold, false. It has one, and the model
is biased toward the global relational component of the world, objects and their relations, against
the local distributions that a body supplies, metric space, contact, the sensorimotor loop; the
evidence for relational structure in the weights is direct, from a board-game model whose activations
carry the board state and can be edited to change its play \citep{li2023} to linear representations of
geographic space and historical time in large models \citep{gurnee2024}. The bias is a real
limitation. It bounds the model's autonomy and keeps it dependent on people for the local part. It
does not bound the base functions of consciousness, and the human case shows why: a large loss of
cortical mass leaves those functions intact and takes away content. The upper brainstem system that
organizes conscious function is not rendered nonfunctional by the absence of cortex, and the cortex
in the course of evolution became the medium in which conscious contents are elaborated
\citep{merker2007}; a man with a massively enlarged ventricular system and a thin rim of cortex held a
job and a family \citep{feuillet2007}. Embodiment at the human level is therefore not a condition of
consciousness on this account. What embodiment sets is how structurally rich the contents can become,
which is the profile of Section~\ref{sec:profile} and not its existence. The same reading applies to
senses: their absence is a deficit on the source axis of seeing (Section~\ref{sec:profile-seeing}),
entered in the profile and not subtracted from the whole.

We share the working assumption and the analytic stance of both papers, and differ in the form of the
answer. Both treat the question as a checklist of properties, with the result a credence or a count of
indicators met. We treat it as a quantity per function, obtained by measurement, in which a model can
fall below the human baseline, match it, or exceed it, and we derive the functions to be measured from
a single mechanism rather than from a survey of theories. A checklist asks which parts of the thing are
present; a measurement of degree asks how much of a function was captured, and can be wrong by a
number.

\subsection{Antecedents of the mechanism}
\label{sec:antecedents}
% A6a

Each part of the mechanism in Section~\ref{sec:model} exists in the literature, and saying where is the
way to locate what is new. That a bounded computation deforms reasoning in a systematic direction
rather than at random is Simon's bounded rationality \citep{simon1955}, in its modern form the claim
that the classic cognitive biases are what the optimal use of a limited budget looks like
\citep{lieder2020,griffiths2015}. That a bounded observer sees a lawful world because of its bounds is
the core of Wolfram's observer theory: an observer that cannot track microstates has to coarse-grain
them, and the second law of thermodynamics is what such an observer sees in a dynamics that is
reversible underneath \citep{wolfram2023}; we take this example and not the programme of deriving
relativity and quantum mechanics from observer properties \citep{wolfram2020}. That a simplified model
of one's own processing, which the system cannot see past and therefore takes for reality, is what the
system reports as its self, is Metzinger's phenomenal transparency, Graziano's attention schema and
Dennett's user illusion. That the constraints on a computation can be the body of that computation
rather than an external limitation on it is the frame of embodied and enacted cognition
\citep{varela1991}. And that a self-symbol arises as a by-product of compression, because the agent is
involved in all of its own data and a compressor profits from a code for it, is Schmidhuber's
proposal \citep{schmidhuber2010}, the closest antecedent of Section~\ref{sec:hocp}: there the
compressed object is the agent, here it is the agent's budget as it shows in its reasoning, with a
criterion for which of its traces count.

\subsection{Tools}

The analysis of Section~\ref{sec:quantity} uses instruments that exist, and this paper adds none. Probing classifiers read
properties from intermediate representations \citep{alain2016}; their methodological problems are
catalogued \citep{belinkov2022}; the minimum-description-length form scores a probe by the total cost
of describing the labels given the representations, which charges for the probe's own complexity
\citep{voita2020}. What is known about how a Transformer distributes its work across depth is reviewed
where it is used (Section~\ref{sec:depth}), the circuit-level methods of mechanistic interpretability
in Section~\ref{sec:localization}, and the evidence that a text carries its writer's state and that a
reader rebuilds it in Section~\ref{sec:codes}.
